% Motivação

O cálculo do poder estatístico é uma etapa fundamental no planejamento de experimentos, pois garante que os recursos sejam utilizados de forma eficiente e que os resultados obtidos sejam confiáveis. Um experimento com baixo poder estatístico pode levar a conclusões equivocadas, seja pela incapacidade de detectar um efeito real (erro tipo II) ou pela superestimação de um efeito observado em amostras pequenas.

Como vimos, pesquisadores normalmente buscam desenhar experimentos que possuam poder estatístico de 80\%, isto é, quando o experimento tem efeito, iremos concluir que esse efeito existe em 80\% das vezes. Este valor depende de parâmetros que estão no controle do pesquisador na fase do desenho do experimento, como o tamanho da amostra e a proporção de tratados, e de parâmetros que não estão ao alcance do pesquisador, como a variância da variável de interesse. Portanto, essa análise prévia ao experimento é de extrema importância para que os resultados encontrados após sua realização sejam críveis.

% Ruído

Ao longo deste guia, destrinchamos o funcionamento da análise de poder para diferentes desenhos de experimentos. A medida em que se aumenta a complexidade dos desenhos de experimento, maior é a presença de ruído, o que diminui o poder do experimento. 

% Fórmula Analítica

Existem duas formas de realizarmos a análise de poder: pela fórmula analítica ou por meio de simulações. Essencialmente, para calcularmos o MDE por meio da fórmula analítica, basta sabermos a fórmula analítica da variância do estimador que estaremos utilizando. Este método é útil para quando não temos acesso a dados antes da realização do nosso experimento, nos permitindo calcular o MDE em termos de desvio padrão.

% Simulação

Quando possuímos acesso aos dados, podemos calcular o MDE através de simulações. Simulando a alocação de tratados e controles no experimento, conseguimos estimar a variância do estimador que iremos utilizar. Esse método também é útil para quando temos fórmulas analíticas da variância do estimador muito complexas.

Ao determinar previamente o poder do experimento, os pesquisadores podem ajustar o tamanho da amostra, considerar a variabilidade dos dados e garantir que os efeitos de interesse sejam detectáveis com um grau adequado de confiança. Isso é particularmente relevante em estudos aplicados, onde recursos financeiros, tempo e disponibilidade de participantes são limitados.

Além disso, o planejamento adequado do poder estatístico melhora a reprodutibilidade dos estudos, reduz o risco de obter resultados inconclusivos e contribui para a credibilidade das análises realizadas. Dessa forma, o cálculo do poder não é apenas uma formalidade estatística, mas uma prática essencial para conduzir experimentos robustos e com impacto significativo.